\begin{titlepage}
\begin{abstract}
%
La creazione di personaggi non giocanti (NPC) credibili nei videogiochi è stata tradizionalmente affrontata tramite macchine a stati ed alberi decisionali, strumenti che generano comportamenti ripetitivi, non si adattano a contesti nuovi e risultano privi di memoria e apprendimento.
Questo lavoro di tesi presenta un'architettura per NPC in ambienti 3D che integra memoria persistente, recupero intelligente dei ricordi, pianificazione gerarchica e riflessione periodica sui propri comportamenti. Il sistema utilizza un modello linguistico(LLM) in esecuzione locale, privo di dipendenze da API cloud, e include un meccanismo di adattamento automatico a scene non preconfigurate.
I risultati mostrano che l'NPC sviluppa comportamenti coerenti con i propri bisogni fisiologici dinamici, mostra capacità di apprendimento tramite le riflessioni e si adatta ad ambienti diversi senza richiedere configurazioni manuali. L'architettura proposta risulta così modulare, estensibile e completamente autonoma da servizi esterni.
%
\\[1cm]
\end{abstract}
\end{titlepage}